\section{Practical Trade-offs and Recommendations}

\subsection{Decision Matrix}

The following table summarizes the key differences across all three platforms
for the chain-of-contracts pattern:

\begin{table}[h]
\centering
\small
\begin{tabular}{|l|c|c|c|}
\hline
\textbf{Factor} & \textbf{Ethereum} & \textbf{Polkadot} & \textbf{TON} \\
\hline
Atomic chain execution & Yes & Yes & No \\
Reentrancy protection & Manual & Built-in (mostly) & Structural \\
Cross-contract complexity & Low & Medium & High \\
Event verification & Trivial & Moderate & Indirect \\
Max sharding & None & Para-chain & Infinite \\
User base & Largest & Growing & Niche \\
Tooling maturity & Mature & Maturing & Early \\
Gas model clarity & Simple & Moderate & Complex \\
Storage cost model & One-time & Deposit & Rent \\
\hline
\end{tabular}
\caption{Comprehensive comparison across platforms}
\end{table}

\subsection{When to Use Each Platform}

\begin{itemize}
    \item \textbf{Ethereum} is the best choice for:
    \begin{itemize}
        \item Rapid prototyping and student education.
        \item Atomic chains where intermediate state must not be persisted on failure.
        \item Applications requiring maximum developer tooling and community support.
        \item DeFi protocols where synchronous composability is essential.
    \end{itemize}

    \item \textbf{Polkadot (ink!)} is the best choice for:
    \begin{itemize}
        \item Projects requiring cross-chain interoperability via XCM.
        \item Applications needing formal verification (Rust's type system helps).
        \item Teams already familiar with Rust.
        \item Scenarios where upgradeable smart contracts (via \texttt{set\_code\_hash}) are needed.
    \end{itemize}

    \item \textbf{TON} is the best choice for:
    \begin{itemize}
        \item High-throughput applications serving millions of users.
        \item Systems where sharding is required today, not as a future roadmap item.
        \item Applications that can tolerate asynchronous message passing.
        \item Projects that benefit from TON's native Telegram integration.
    \end{itemize}
\end{itemize}

\subsection{Hybrid Approaches}

For production systems, platforms can be combined:

\begin{itemize}
    \item \textbf{ETH + TON bridges}: Use Ethereum for the atomic chain core
          and TON for high-volume message relaying to end users.
    \item \textbf{Polkadot + TON}: Deploy the chain-of-contracts as a parachain
          (synchronous, sovereign) and bridge messages to TON for distribution.
    \item \textbf{TON with bounce guards}: Implement the A $\to$ B $\to$ C chain
          in TON using bounce messages and confirmation timeouts, accepting
          eventual consistency for scalability.
\end{itemize}

\subsection{Educational Recommendations}

For a chain-of-contracts lab, we recommend:

\begin{enumerate}
    \item \textbf{Implement on Ethereum first} -- the synchronous model is
          easier to understand and debug. Students can verify the chain with
          a single transaction receipt.
    \item \textbf{Port to Polkadot} -- demonstrates how the same pattern
          translates to a different VM with stronger type safety and a
          different cross-contract API.
    \item \textbf{Analyze TON theoretically (this document)} -- understand
          how the pattern must be redesigned for the actor model, without
          requiring FunC proficiency.
\end{enumerate}

This three-step approach gives students both practical implementation skills
and a theoretical framework for evaluating smart contract platforms.